\begin{afterword}
\summaryhead

The post returns to reductionism, which it restates as the thesis that we use many-levelled models for reasons of computation while reality has a single level. It asks whether, when you pick up a cup, it is your hand or your fingers that pick it up, and jokes that scientists have found the hand's force to be zero, since the fingers, thumb and palm account for all of it.

The author's point is that once you can see how a higher level reduces to a lower one, as you can see a hand as fingers, thumb and palm, the question becomes silly. A reasoner who only knows the fact of reduction, from separate scanners or from being told, can still imagine the hand moving off alone. That is a fact about the reasoner's map. What seems possible to a particular person is not what is logically or physically possible, just as 235757 may seem possibly prime to someone who has not factored it. The post ends by parodying philosophers who deny that hands exist, or who posit contingent laws by which fingers bring a hand into being, and calls both the Mind Projection Fallacy.

\responsehead

The post's central point is sound and familiar. What someone can imagine depends on what they know, so the imaginability of the hand without the fingers shows nothing about hands. Putnam wrote in 1975 that ``Conceivability is no proof of logical possibility,'' and the author quotes that sentence two days later. The number example is correct, and the joke about the hand's missing force has the form of a real argument, the exclusion problem (how a whole can cause anything its parts do not already cause), to which the post gives one of the standard answers: if the hand is its parts, nothing is excluded.

The difficulty is the choice of example. In the case of the hand, no one disputes the post. The two positions it parodies are real only for the mind. Eliminative materialists deny that beliefs and desires exist because they expect those states not to reduce; they do not deny things whose reduction is plain. Dualists who posit laws linking experience to physical states hold them for consciousness only, and Chalmers writes that ``it seems that we cannot coherently conceive'' of a world physically like ours without water. Applied to hands, both positions lose the reasons their real holders give, and the verdict that they take intuitions for facts is reached only on those versions. Whether the mind is like the hand is the question in dispute, and this post does not argue it; the book returns to it in the zombie posts.

One claim is stronger than the argument needs. The post says that the levels are bound by ``physical identity'' and that ``hand'' and ``fingers and thumb and palm'' differ only in point of view. That is the metaphysical thesis of composition as identity, which the \textsc{Stanford Encyclopedia} calls ``by no means an uncontroversial view.'' The post's reasoning needs only that the hand is made of its parts.

In short: a sound point about imagination and possibility, made on the hand, where it is not disputed, and used against parodies of views that their real holders apply only to the mind.
\end{afterword}
